\documentclass{article}
\usepackage[a4paper,margin=0.75in]{geometry}
\usepackage{amsmath}
\usepackage{array}
\usepackage{correlator-diagrams}

\setlength{\parindent}{0pt}
\renewcommand{\arraystretch}{1.8}

\newcommand{\FeatureThreePoint}[1]{%
  \ThreePointCorr[
    #1,
    line=ewboson,
    leg-labels={W_\mu,Z_\nu,W_\rho},
    indices={\alpha,\beta,\gamma},
    momentum-labels={q_1,q_2,q_3},
    three-point-momentum-flow=all-in,
    three-point-momentum-layout=outward,
    leg-styles={{epspol,fieldA},ewboson,ewboson},
    three-point-external-length=1.08,
    propagator-arrow-sizes={13pt,,}
  ]%
}

\newcommand{\RawEndcappSample}[1][]{%
  \begin{tikzpicture}[baseline=(current bounding box.center),scale=1]
    \begin{feynhand}
      \vertex (a) at (0,0) {};
      \vertex (b) at (1.85,0) {};
      \propag[ewboson,endcapp,#1] (a) to (b);
      \vertex[dot] at (a) {};
      \vertex[dot] at (b) {};
    \end{feynhand}
  \end{tikzpicture}%
}

\begin{document}

\section*{Three-Point Feature Gallery}

All four orientations, with the special slot-1 leg styled as
\texttt{epspol}. Slot labels, indices, momenta, and leg styles remain attached
to slots 1, 2, and 3 after rotation.

\[
\begin{array}{cc}
  \text{right (compatibility layout)} &
  \text{left} \\
  \FeatureThreePoint{three-point-orientation=right} &
  \FeatureThreePoint{three-point-orientation=left} \\[2.5em]
  \text{up} &
  \text{down} \\
  \FeatureThreePoint{three-point-orientation=up} &
  \FeatureThreePoint{three-point-orientation=down}
\end{array}
\]

\newpage

Scaling and external-leg length. The first row scales the full drawing; the
second row changes only the special slot-1 external leg.

\[
\begin{array}{cc}
  \text{scale 0.72} &
  \text{scale 1.24} \\
  \FeatureThreePoint{scale=0.72,three-point-orientation=right} &
  \FeatureThreePoint{scale=1.24,three-point-orientation=right} \\[2.5em]
  \text{short slot 1} &
  \text{longer slot 1} \\
  \FeatureThreePoint{three-point-external-length=0.78} &
  \FeatureThreePoint{three-point-external-length=1.45}
\end{array}
\]

\newpage

Double-endcap tuning on three-point vertices. The inner cap is intentionally
longer than the terminal cap, and the gap between the bars stays blank.

\[
\begin{array}{cc}
  \text{compact caps} &
  \text{longer caps and wider gap} \\
  \FeatureThreePoint{endpoint-cap-length=0.48,endpoint-double-cap-inner-length=0.72,endpoint-double-cap-gap=2.0pt} &
  \FeatureThreePoint{endpoint-cap-length=0.86,endpoint-double-cap-inner-length=1.24,endpoint-double-cap-gap=5.0pt} \\[2.5em]
  \text{rotated compact caps} &
  \text{rotated wider gap} \\
  \FeatureThreePoint{three-point-orientation=up,endpoint-cap-length=0.48,endpoint-double-cap-inner-length=0.72,endpoint-double-cap-gap=2.0pt} &
  \FeatureThreePoint{three-point-orientation=down,endpoint-cap-length=0.86,endpoint-double-cap-inner-length=1.24,endpoint-double-cap-gap=5.0pt}
\end{array}
\]

\newpage

Raw \texttt{endcapp} propagators also reserve blank space between the two
bars, independent of the topology-managed drawing path.

\[
\begin{array}{cccc}
  \text{default} &
  \text{color before} &
  \text{color after} &
  \text{wide gap} \\
  \RawEndcappSample &
  \RawEndcappSample[fieldA] &
  \RawEndcappSample[color=blue!65!black] &
  \RawEndcappSample[
    endpoint cap length=0.78,
    endpoint double cap inner length=1.12,
    endpoint double cap gap=5.2pt
  ]
\end{array}
\]

\end{document}
